% EGE 2027, task 17. Source IDs and URLs are kept in comments.
% Editable task bank. Append new BankTask blocks at the end of the document.
\documentclass[12pt,a4paper]{article}
\usepackage{physics-textbook}
\usepackage{graphicx,needspace}
\newif\ifShowSolutions
\ifdefined\TutorEdition
  \ShowSolutionstrue
\else
  \ShowSolutionsfalse
\fi
\newcounter{banktask}
\newcommand{\BankTask}[3]{%
  \refstepcounter{banktask}%
  \Needspace{35mm}\par\medskip\noindent\begin{minipage}{\linewidth}%
  \noindent\textbf{\thebanktask.}\quad #2\par
  \ifShowSolutions
    #3\par
    \ifdefined\TutorEdition
      \ifcsname TutorSource#1\endcsname
        \smallskip{\small\csname TutorSource#1\endcsname\par}%
      \fi
    \fi
  \fi
  \end{minipage}\par\medskip
}
\begin{document}


% Source 7D7448: https://ege.fipi.ru/bank/questions.php?search=1&proj=BA1F39653304A5B041B656915DC36B38&qid=7D7448
\BankTask{7D7448}{
Установите соответствие между физическими величинами и формулами, по которым их можно рассчитать ($E$ — энергия фотона; $h$ — постоянная Планка; $p$ — импульс фотона).

\TaskMatchingColumns[.53\linewidth]{Физические величины}{
А) длина волны фотона\par
Б) частота фотона
}{Формулы}{
1) $\dfrac{p}{h}$\par
2) $\dfrac{h}{p}$\par
3) $\dfrac{E}{p}$\par
4) $\dfrac{E}{h}$
}

\medskip
\begin{tabular}{|c|c|}
\hline
А & Б \\
\hline
\phantom{2} & \phantom{4} \\
\hline
\end{tabular}
}{
\smallskip\noindent\textbf{Решение от GPT}\par
Из соотношения $p=h/\lambda$ получаем $\lambda=h/p$, поэтому А соответствует формуле 2. По формуле Планка $E=h\nu$, откуда $\nu=E/h$, поэтому Б соответствует формуле 4. Ответ: $24$.
}

% Source 64FC99: https://ege.fipi.ru/bank/questions.php?search=1&proj=BA1F39653304A5B041B656915DC36B38&qid=64FC99
\BankTask{64FC99}{
Установите соответствие между физическими величинами и формулами, по которым их можно рассчитать ($\lambda$ — длина волны фотона; $c$ — скорость света в вакууме; $h$ — постоянная Планка).

\TaskMatchingColumns[.53\linewidth]{Физические величины}{
А) частота фотона\par
Б) импульс фотона
}{Формулы}{
1) $\dfrac{h}{\lambda}$\par
2) $\dfrac{h\lambda}{c}$\par
3) $\dfrac{c}{\lambda}$\par
4) $c\lambda$
}

\medskip
\begin{tabular}{|c|c|}
\hline
А & Б \\
\hline
\phantom{3} & \phantom{1} \\
\hline
\end{tabular}
}{
\smallskip\noindent\textbf{Решение от GPT}\par
Из $c=\lambda\nu$ следует $\nu=c/\lambda$, поэтому А соответствует формуле 3. Импульс фотона $p=h/\lambda$, поэтому Б соответствует формуле 1. Ответ: $31$.
}

% Source D4FFF9: https://ege.fipi.ru/bank/questions.php?search=1&proj=BA1F39653304A5B041B656915DC36B38&qid=D4FFF9
\BankTask{D4FFF9}{
Установите соответствие между графиками, представленными на рисунках, и зависимостями, которые они могут выражать.

\par\medskip\TaskMatchHeading{Графики}

А) рисунок 1

Б) рисунок 2

\par\medskip\TaskMatchHeading{Зависимости}

1) зависимость силы фототока от напряжения между электродами при неизменной освещённости

2) зависимость энергии фотона от длины волны

3) зависимость максимальной энергии фотоэлектронов от частоты света

4) зависимость энергии фотона от частоты света

\begin{tabular}{|c|c|}
\hline
А & Б \\
\hline
\phantom{2} & \phantom{3} \\
\hline
\end{tabular}

\medskip
\begin{center}
\begin{tikzpicture}[book diagram, x=0.024cm, y=0.024cm]
\draw[-{book tip}] (0,0) -- (150,0);
\draw[-{book tip}] (0,0) -- (0,127);
\node[below left] at (0,0) {$0$};
\draw[line width=0.5pt] (21,125) .. controls (22,118) and (21,105) .. (23,94) .. controls (26,73) and (32,58) .. (42,48) .. controls (54,36) and (69,29) .. (87,25) .. controls (104,21) and (125,19) .. (143,19);
\draw (0,103) -- (22,103);
\draw (22,103) -- (22,80);
\draw (22,32) -- (22,8);
\draw (0,23) -- (7,23);
\draw (32,23) -- (55,23);
\draw (80,23) -- (104,23);
\draw (104,23) -- (104,0);
\node[below] at (72,-18) {А) Рисунок 1};
\end{tikzpicture}
\hspace{1cm}
\begin{tikzpicture}[book diagram, x=0.024cm, y=0.024cm]
\draw[-{book tip}] (0,0) -- (163,0);
\draw[-{book tip}] (0,0) -- (0,127);
\node[below left] at (0,0) {$0$};
\draw[line width=0.7pt] (0,0) -- (52,0) -- (151,99);
\node[below] at (76,-18) {Б) Рисунок 2};
\end{tikzpicture}
\end{center}
}{
\smallskip\noindent\textbf{Решение от GPT}\par
График А — убывающая гипербола: $E=hc/\lambda$, поэтому ему соответствует зависимость 2. График Б описывает линейный рост максимальной энергии фотоэлектронов выше пороговой частоты: $K_{\max}=h\nu-A_{\text{вых}}$, то есть зависимость 3. Зависимость 1 представляет рост фототока с выходом на насыщение, а зависимость 4 — прямую $E=h\nu$, проходящую через начало координат; они не соответствуют данным графикам. Ответ: $23$.
}

% Source 369826: https://ege.fipi.ru/bank/questions.php?search=1&proj=BA1F39653304A5B041B656915DC36B38&qid=369826
\BankTask{369826}{
На металлическую пластинку падает пучок монохроматического света. При этом наблюдается явление фотоэффекта.

На графике А представлена зависимость энергии фотонов, падающих на катод, от физической величины \(x_1\), а на графике Б — зависимость максимальной кинетической энергии фотоэлектронов от физической величины \(x_2\).

Какая из физических величин отложена на горизонтальной оси на графике А и какая — на графике Б?

\TaskMatchingColumns[.53\linewidth]{Графики}{
А) график А\par
Б) график Б
}{Физическая величина \(x\)}{
1) частота\par
2) длина волны\par
3) массовое число\par
4) заряд ядра
}

\medskip
\begin{tabular}{|c|c|}
\hline
А & Б \\
\hline
\phantom{2} & \phantom{1} \\
\hline
\end{tabular}

\medskip
\begin{tikzpicture}[book diagram,x=0.45cm,y=0.23cm]
\node at (3.5,20.5) {А)};
\draw[-{book tip}] (0,0) -- (7.3,0) node[below] {$x_1$};
\draw[-{book tip}] (0,0) -- (0,18.5) node[left] {$E$};
\node[below left] at (0,0) {$0$};
\foreach \y in {0.5,1,...,17.5} {\draw (-0.09,\y) -- (0.09,\y);}
\foreach \y in {5,10,15} {\node[left] at (-0.14,\y) {$\y$};}
\foreach \xx/\yy in {1/15,2/7.5,3/5,4/3.75,5/3,6/2.5} {
\draw[thin] (0,\yy) -- (\xx,\yy) -- (\xx,0);
}
\draw[thick] plot[domain=0.82:7,samples=100] (\x,{15/\x});
\foreach \x/\y in {1/15,2/7.5,3/5,4/3.75,5/3,6/2.5}{\fill (\x,\y) circle[radius=1pt];}
\end{tikzpicture}
\qquad
\begin{tikzpicture}[book diagram,x=0.45cm,y=0.45cm]
\node at (3.5,5.25) {Б)};
\draw[-{book tip}] (0,0) -- (7.1,0) node[below] {$x_2$};
\draw[-{book tip}] (0,-1.25) -- (0,4.7) node[left] {$E$};
\node[left] at (0,0) {$0$};
\draw[thick] (0,0) -- (2,0) -- (5.7,3.7);
\end{tikzpicture}
}{
\smallskip\noindent\textbf{Решение от GPT}\par
Энергия фотона \(E=h\nu=hc/\lambda\), поэтому убывающая гипербола А соответствует длине волны (2). По уравнению Эйнштейна \(K_{\max}=h\nu-A_{\text{вых}}\), поэтому возрастающая прямая Б соответствует частоте (1). Массовое число (3) и заряд ядра (4) не являются переменными в этих зависимостях. Ответ: А — 2, Б — 1.
}

% Source BD7F43: https://ege.fipi.ru/bank/questions.php?search=1&proj=BA1F39653304A5B041B656915DC36B38&qid=BD7F43
\BankTask{BD7F43}{
Установите соответствие между видами радиоактивного распада и уравнениями, описывающими этот процесс.

\TaskMatchingColumns[.25\linewidth]{Виды распада}{
А) альфа-распад\par
Б) бета-распад
}{Уравнения}{
1) \({}^{238}_{92}\mathrm{U}+{}^{22}_{10}\mathrm{Ne}\to{}^{256}_{102}\mathrm{No}+4{}^{1}_{0}\mathrm{n}\)\par
2) \({}^{238}_{93}\mathrm{Np}\to{}^{238}_{94}\mathrm{Pu}+{}^{0}_{-1}\mathrm{e}+\widetilde{\nu}_e\)\par
3) \({}^{227}_{89}\mathrm{Ac}\to{}^{223}_{87}\mathrm{Fr}+{}^{4}_{2}\mathrm{He}\)\par
4) \({}^{209}_{83}\mathrm{Bi}+{}^{2}_{1}\mathrm{H}\to{}^{105}_{43}\mathrm{Tc}+{}^{102}_{41}\mathrm{Nb}+4{}^{1}_{0}\mathrm{n}\)
}

\medskip
\begin{tabular}{|c|c|}
\hline
А & Б \\
\hline
\phantom{3} & \phantom{2} \\
\hline
\end{tabular}
}{
\smallskip\noindent\textbf{Решение от GPT}\par
При альфа-распаде испускается ядро \({}^{4}_{2}\mathrm{He}\): массовое число уменьшается на 4, зарядовое — на 2; это уравнение 3. При электронном бета-распаде испускаются электрон и электронное антинейтрино, массовое число сохраняется, а зарядовое увеличивается на 1; это уравнение 2. Уравнения 1 и 4 описывают вынужденные ядерные реакции с налетающей частицей. Ответ: А — 3, Б — 2.
}

% Source FF6C82: https://ege.fipi.ru/bank/questions.php?search=1&proj=BA1F39653304A5B041B656915DC36B38&qid=FF6C82
\BankTask{FF6C82}{
Выберите среди приведённых во втором столбце ядерных реакций те, которые являются примерами реакций альфа- и бета-распада.

\TaskMatchingColumns[.25\linewidth]{Вид ядерной реакции}{
А) альфа-распад\par
Б) бета-распад
}{Ядерные реакции}{
1) ${}_{77}^{176}\mathrm{Ir}\to{}_{75}^{172}\mathrm{Re}+{}_{2}^{4}\mathrm{He}$\par
2) ${}_{92}^{238}\mathrm{U}+{}_{10}^{22}\mathrm{Ne}\to{}_{102}^{256}\mathrm{No}+4{}_{0}^{1}n$\par
3) ${}_{48}^{113}\mathrm{Cd}+{}_{0}^{1}n\to{}_{48}^{114}\mathrm{Cd}+\gamma$\par
4) ${}_{71}^{178}\mathrm{Lu}\to{}_{72}^{178}\mathrm{Hf}+{}_{-1}^{0}e+\tilde{\nu}_e$
}

\begin{tabular}{|c|c|}
\hline
А & Б \\
\hline
\phantom{0} & \phantom{0} \\
\hline
\end{tabular}
}{
\smallskip\noindent\textbf{Решение от GPT}\par
Реакция 1 является альфа-распадом: испускается ${}_{2}^{4}\mathrm{He}$, поэтому $A$ уменьшается на 4, а $Z$ — на 2. Реакция 4 является электронным бета-распадом: испускаются электрон и антинейтрино, $A$ сохраняется, а $Z$ увеличивается на 1. Реакция 2 — слияние ядер с испусканием нейтронов, реакция 3 — захват нейтрона с испусканием гамма-кванта. Ответ: А — 1, Б — 4.
}

% Source F926FD: https://ege.fipi.ru/bank/questions.php?search=1&proj=BA1F39653304A5B041B656915DC36B38&qid=F926FD
\BankTask{F926FD}{
На рисунке изображена упрощённая диаграмма нижних энергетических уровней атома. Стрелками отмечены некоторые возможные переходы атома между этими уровнями.

Установите соответствие между процессами поглощения света наименьшей длины волны и излучения света наименьшей частоты и энергией соответствующего фотона.

\TaskMatchingColumns[.53\linewidth]{Процесс}{
А) поглощение света наименьшей длины волны\par
Б) излучение света наименьшей частоты
}{Энергия фотона}{
1) $E_1-E_0$\par
2) $E_2-E_0$\par
3) $E_3-E_0$\par
4) $E_4-E_0$
}

\begin{tabular}{|c|c|}
\hline
А & Б \\
\hline
\phantom{0} & \phantom{0} \\
\hline
\end{tabular}

\begin{center}
\begin{tikzpicture}[book diagram,x=0.035cm,y=0.035cm]
\draw (27,0)--(27,175);
\draw[dash pattern=on 24pt off 24pt] (32,159)--(144,159);
\node[left] at (23,159) {$0$};
\draw (27,137)--(144,137);
\draw (27,121)--(144,121);
\draw (27,101)--(144,101);
\draw (27,67)--(144,67);
\draw (27,17)--(144,17);
\node[left] at (23,137) {$E_4$};
\node[left] at (23,121) {$E_3$};
\node[left] at (23,101) {$E_2$};
\node[left] at (23,67) {$E_1$};
\node[left] at (23,17) {$E_0$};
\draw[-{book tip}] (52,17)--(52,137);
\draw[-{book tip}] (75,17)--(75,101);
\draw[-{book tip}] (97,67)--(97,17);
\draw[-{book tip}] (118,121)--(118,17);
\end{tikzpicture}
\end{center}
}{
\smallskip\noindent\textbf{Решение от GPT}\par
Энергия фотона равна $\varepsilon=h\nu=hc/\lambda=\lvert E_f-E_i\rvert$. Среди показанных переходов поглощения $E_0\to E_4$ и $E_0\to E_2$ наименьшей длине волны соответствует наибольшая энергия $E_4-E_0$ (вариант 4, а не 2). Среди переходов излучения $E_1\to E_0$ и $E_3\to E_0$ наименьшей частоте соответствует меньшая энергия $E_1-E_0$ (вариант 1, а не 3). Ответ: А — 4, Б — 1.
}

% Source AB5BFF: https://ege.fipi.ru/bank/questions.php?search=1&proj=BA1F39653304A5B041B656915DC36B38&qid=AB5BFF
\BankTask{AB5BFF}{
На рисунке изображена упрощённая диаграмма нижних энергетических уровней атома. Нумерованными стрелками отмечены некоторые возможные переходы атома между этими уровнями. Какой из этих четырёх переходов связан с поглощением света наименьшей частоты, а какой — с излучением света наибольшей частоты?

Установите соответствие между процессами поглощения и испускания света и стрелками, указывающими энергетические переходы атома.

\begin{center}
\TaskMatchingColumns[0.58\linewidth]{Процесс}{
А) поглощение света наименьшей частоты\par
Б) излучение света наибольшей частоты
}{Энергетический переход}{
1) 1\par
2) 2\par
3) 3\par
4) 4
}

\medskip
\begin{tabular}{|c|c|}
\hline
А & Б \\
\hline
\phantom{3} & \phantom{4} \\
\hline
\end{tabular}
\end{center}

\begin{center}
\begin{tikzpicture}[book diagram,x=0.025cm,y=0.025cm]
\draw (0,-9) -- (0,203);
\draw[dashed] (0,189) -- (125,189);
\node[left] at (0,189) {$0$};
\draw (0,157) -- (125,157);
\node[left] at (0,157) {$E_4$};
\draw (0,141) -- (125,141);
\node[left] at (0,136) {$E_3$};
\draw (0,110) -- (125,110);
\node[left] at (0,105) {$E_2$};
\draw (0,63) -- (125,63);
\node[left] at (0,63) {$E_1$};
\draw (0,0) -- (125,0);
\node[left] at (0,0) {$E_0$};
\draw[-{book tip}] (26,0) -- (26,141);
\draw[-{book tip}] (51,110) -- (51,0);
\draw[-{book tip}] (79,0) -- (79,63);
\draw[-{book tip}] (105,157) -- (105,0);
\node[below] at (26,0) {$1$};
\node[below] at (51,0) {$2$};
\node[below] at (79,0) {$3$};
\node[below] at (105,0) {$4$};
\end{tikzpicture}
\end{center}
}{
\smallskip\noindent\textbf{Решение от GPT}\par
Частота света определяется разностью энергий уровней: $\nu=|E_f-E_i|/h$. Поглощению соответствуют переходы 1 ($E_0\to E_3$) и 3 ($E_0\to E_1$), поэтому наименьшая частота соответствует переходу 3. Излучению соответствуют переходы 2 ($E_2\to E_0$) и 4 ($E_4\to E_0$), поэтому наибольшая частота соответствует переходу 4. Ответ: А — 3, Б — 4.
}

% Source 6BBB75: https://ege.fipi.ru/bank/questions.php?search=1&proj=BA1F39653304A5B041B656915DC36B38&qid=6BBB75
\BankTask{6BBB75}{
На рисунке изображена упрощённая диаграмма нижних энергетических уровней атома. Нумерованными стрелками отмечены некоторые возможные переходы атома между этими уровнями. Какие из этих четырёх переходов связаны с излучением света с наибольшей длиной волны и поглощением света с наименьшей энергией?

Установите соответствие между процессами поглощения и излучения света и энергетическими переходами атома, указанными стрелками.

\begin{center}
\TaskMatchingColumns[0.58\linewidth]{Процессы}{
А) излучение света с наибольшей длиной волны\par
Б) поглощение света с наименьшей энергией
}{Энергетические переходы}{
1) 1\par
2) 2\par
3) 3\par
4) 4
}

\medskip
\begin{tabular}{|c|c|}
\hline
А & Б \\
\hline
\phantom{1} & \phantom{3} \\
\hline
\end{tabular}
\end{center}

\begin{center}
\begin{tikzpicture}[book diagram,x=0.025cm,y=0.025cm]
\draw (0,-13) -- (0,192);
\draw[dashed] (0,177) -- (117,177);
\node[left] at (0,177) {$0$};
\draw (0,148) -- (117,148);
\node[left] at (0,148) {$E_4$};
\draw (0,133) -- (117,133);
\node[left] at (0,128) {$E_3$};
\draw (0,104) -- (117,104);
\node[left] at (0,99) {$E_2$};
\draw (0,59) -- (117,59);
\node[left] at (0,59) {$E_1$};
\draw (0,0) -- (117,0);
\node[left] at (0,0) {$E_0$};
\draw[-{book tip}] (25,59) -- (25,0);
\draw[-{book tip}] (49,104) -- (49,0);
\draw[-{book tip}] (73,0) -- (73,133);
\draw[-{book tip}] (97,0) -- (97,148);
\node[below] at (25,0) {$1$};
\node[below] at (49,0) {$2$};
\node[below] at (73,0) {$3$};
\node[below] at (97,0) {$4$};
\end{tikzpicture}
\end{center}
}{
\smallskip\noindent\textbf{Решение от GPT}\par
Энергия фотона равна $\varepsilon=|E_f-E_i|$, а длина волны $\lambda=hc/\varepsilon$. Излучению соответствуют переходы 1 ($E_1\to E_0$) и 2 ($E_2\to E_0$); меньшая энергия и наибольшая длина волны соответствуют переходу 1. Поглощению соответствуют переходы 3 ($E_0\to E_3$) и 4 ($E_0\to E_4$); наименьшая энергия соответствует переходу 3. Ответ: А — 1, Б — 3.
}

% Source BE27DA: https://ege.fipi.ru/bank/questions.php?search=1&proj=BA1F39653304A5B041B656915DC36B38&qid=BE27DA
\BankTask{BE27DA}{
На рисунке изображена упрощённая диаграмма нижних энергетических уровней атома. Нумерованными стрелками отмечены некоторые возможные переходы атома между этими уровнями. Какой из этих четырёх переходов связан с поглощением кванта света наибольшей частоты, а какой — с излучением кванта света наименьшей частоты?

Установите соответствие между процессами поглощения и излучения света и энергетическими переходами атома, указанными стрелками.

\begin{center}
\TaskMatchingColumns[0.58\linewidth]{Процессы}{
А) поглощение кванта света наибольшей частоты\par
Б) излучение кванта света наименьшей частоты
}{Энергетические переходы}{
1) 1\par
2) 2\par
3) 3\par
4) 4
}

\medskip
\begin{tabular}{|c|c|}
\hline
А & Б \\
\hline
\phantom{4} & \phantom{1} \\
\hline
\end{tabular}
\end{center}

\begin{center}
\begin{tikzpicture}[book diagram,x=0.025cm,y=0.025cm]
\draw (0,-13) -- (0,184);
\draw[dashed] (0,170) -- (112,170);
\node[left] at (0,170) {$0$};
\draw (0,142) -- (112,142);
\node[left] at (0,142) {$E_4$};
\draw (0,127) -- (112,127);
\node[left] at (0,122) {$E_3$};
\draw (0,99) -- (112,99);
\node[left] at (0,94) {$E_2$};
\draw (0,57) -- (112,57);
\node[left] at (0,57) {$E_1$};
\draw (0,0) -- (112,0);
\node[left] at (0,0) {$E_0$};
\draw[-{book tip}] (23,57) -- (23,0);
\draw[-{book tip}] (46,99) -- (46,0);
\draw[-{book tip}] (71,0) -- (71,127);
\draw[-{book tip}] (93,0) -- (93,142);
\node[below] at (23,0) {$1$};
\node[below] at (46,0) {$2$};
\node[below] at (71,0) {$3$};
\node[below] at (93,0) {$4$};
\end{tikzpicture}
\end{center}
}{
\smallskip\noindent\textbf{Решение от GPT}\par
Частота кванта света равна $\nu=|E_f-E_i|/h$. Поглощению соответствуют переходы 3 ($E_0\to E_3$) и 4 ($E_0\to E_4$), поэтому наибольшая частота соответствует переходу 4. Излучению соответствуют переходы 1 ($E_1\to E_0$) и 2 ($E_2\to E_0$), поэтому наименьшая частота соответствует переходу 1. Ответ: А — 4, Б — 1.
}

% Source 4Ae2Ae: https://ege.fipi.ru/bank/questions.php?search=1&proj=BA1F39653304A5B041B656915DC36B38&qid=4Ae2Ae
\BankTask{4Ae2Ae}{
На рисунке изображена упрощённая диаграмма нижних энергетических уровней атома. Стрелками отмечены некоторые возможные переходы атома между этими уровнями.

Установите соответствие между процессами поглощения света наибольшей длины волны и излучения света с наименьшей энергией и энергией соответствующего фотона.

\begin{center}
\TaskMatchingColumns[0.58\linewidth]{Процесс}{
А) поглощение света наибольшей длины волны\par
Б) излучение света с наименьшей энергией
}{Энергия фотона}{
1) $E_1-E_0$\par
2) $E_2-E_0$\par
3) $E_3-E_0$\par
4) $E_4-E_0$
}
\end{center}

\begin{center}
\begin{tabular}{|c|c|}
\hline
А & Б \\
\hline
\phantom{2} & \phantom{1} \\
\hline
\end{tabular}
\end{center}

\begin{center}
\begin{tikzpicture}[book diagram, x=1cm, y=1cm]
\draw[thin] (0,-0.4) -- (0,2.55);
\draw[thin] (-0.4,0) -- (1.8,0);
\draw[dashed] (0,2.42) -- (1.8,2.42);
\draw (0,2) -- (1.55,2);
\draw (0,1.72) -- (1.55,1.72);
\draw (0,1.23) -- (1.55,1.23);
\draw (0,0.70) -- (1.55,0.70);
\node[left] at (-0.05,2.34) {$0$};
\node[left] at (-0.05,2) {$E_4$};
\node[left] at (-0.05,1.72) {$E_3$};
\node[left] at (-0.05,1.23) {$E_2$};
\node[left] at (-0.05,0.70) {$E_1$};
\node[below left] at (-0.05,0) {$E_0$};
\draw[-{book tip}] (0.40,0.70) -- (0.40,0);
\draw[-{book tip}] (0.74,2) -- (0.74,0);
\draw[-{book tip}] (0.98,0) -- (0.98,1.23);
\draw[-{book tip}] (1.30,0) -- (1.30,1.72);
\end{tikzpicture}
\end{center}
}{
\smallskip\noindent\textbf{Решение от GPT}\par
Энергия фотона равна разности энергий уровней: $\varepsilon=hc/\lambda$. Среди показанных переходов поглощения $E_0\to E_2$ и $E_0\to E_3$ наибольшей длине волны соответствует меньшая энергия $E_2-E_0$ — вариант 2. Среди переходов излучения $E_1\to E_0$ и $E_4\to E_0$ наименьшая энергия равна $E_1-E_0$ — вариант 1; варианты 3 и 4 соответствуют более высоким энергиям. Ответ: А — 2, Б — 1.
}

% Source Be82c1: https://ege.fipi.ru/bank/questions.php?search=1&proj=BA1F39653304A5B041B656915DC36B38&qid=Be82c1
\BankTask{Be82c1}{
На рисунке изображена упрощённая диаграмма нижних энергетических уровней атома. Стрелками отмечены некоторые возможные переходы атома между этими уровнями.

Установите соответствие между процессами поглощения света наибольшей длины волны и испускания света наименьшей частоты и энергией соответствующего фотона.

\begin{center}
\TaskMatchingColumns[0.58\linewidth]{Процесс}{
А) поглощение света наибольшей длины волны\par
Б) испускание света наименьшей частоты
}{Энергия фотона}{
1) $E_1-E_0$\par
2) $E_2-E_0$\par
3) $E_3-E_0$\par
4) $E_4-E_0$
}
\end{center}

\begin{center}
\begin{tabular}{|c|c|}
\hline
А & Б \\
\hline
\phantom{2} & \phantom{1} \\
\hline
\end{tabular}
\end{center}

\begin{center}
\begin{tikzpicture}[book diagram, x=1cm, y=1cm]
\draw[thin] (0,-0.4) -- (0,2.74);
\draw[thin] (-0.4,0) -- (1.92,0);
\draw (0,2.10) -- (1.62,2.10);
\draw (0,1.76) -- (1.62,1.76);
\draw (0,1.34) -- (1.62,1.34);
\draw (0,0.76) -- (1.62,0.76);
\node[left] at (-0.05,2.25) {$E_4$};
\node[left] at (-0.05,1.91) {$E_3$};
\node[left] at (-0.05,1.39) {$E_2$};
\node[left] at (-0.05,0.81) {$E_1$};
\node[below left] at (-0.05,0) {$E_0$};
\draw[-{book tip}] (0.32,0.76) -- (0.32,0);
\draw[-{book tip}] (0.66,1.76) -- (0.66,0);
\draw[-{book tip}] (0.96,0) -- (0.96,1.34);
\draw[-{book tip}] (1.34,0) -- (1.34,2.10);
\end{tikzpicture}
\end{center}
}{
\smallskip\noindent\textbf{Решение от GPT}\par
Используем $\varepsilon=h\nu=hc/\lambda$. Из показанных переходов поглощения $E_0\to E_2$ и $E_0\to E_4$ наибольшая длина волны соответствует энергии $E_2-E_0$ — варианту 2, а вариант 4 даёт меньшую длину волны. Из переходов испускания $E_1\to E_0$ и $E_3\to E_0$ наименьшая частота соответствует энергии $E_1-E_0$ — варианту 1, а вариант 3 даёт большую частоту. Ответ: А — 2, Б — 1.
}

% Source 41596C: https://ege.fipi.ru/bank/questions.php?search=1&proj=BA1F39653304A5B041B656915DC36B38&qid=41596C
\BankTask{41596C}{
На рисунке изображена упрощённая диаграмма нижних энергетических уровней атома. Нумерованными стрелками отмечены некоторые возможные переходы атома между уровнями. Какие из этих четырёх переходов связаны с поглощением света с наименьшей частотой и излучением света с наименьшей энергией фотонов?

Установите соответствие между процессами поглощения и излучения света и энергетическими переходами атома, указанными стрелками.

\begin{center}
\TaskMatchingColumns[0.56\linewidth]{Процессы}{
А) поглощение света с наименьшей частотой\par
Б) излучение света с наименьшей энергией фотонов
}{Энергетические переходы}{
1) 1\par
2) 2\par
3) 3\par
4) 4
}
\end{center}

\begin{center}
\begin{tabular}{|c|c|}
\hline
А & Б \\
\hline
\phantom{1} & \phantom{3} \\
\hline
\end{tabular}
\end{center}

\begin{center}
\begin{tikzpicture}[book diagram, x=1cm, y=1cm]
\draw[thin] (0,-0.36) -- (0,3.54);
\draw (0,0) -- (2.16,0);
\draw (0,1.08) -- (2.16,1.08);
\draw (0,1.90) -- (2.16,1.90);
\draw (0,2.46) -- (2.16,2.46);
\draw (0,2.72) -- (2.16,2.72);
\draw[dashed] (0,3.26) -- (2.16,3.26);
\node[left] at (-0.05,3.26) {$0$};
\node[left] at (-0.05,2.80) {$E_4$};
\node[left] at (-0.05,2.46) {$E_3$};
\node[left] at (-0.05,1.90) {$E_2$};
\node[left] at (-0.05,1.08) {$E_1$};
\node[left] at (-0.05,0) {$E_0$};
\draw[-{book tip}] (0.44,0) -- (0.44,1.08);
\draw[-{book tip}] (0.88,0) -- (0.88,1.90);
\draw[-{book tip}] (1.34,2.46) -- (1.34,0);
\draw[-{book tip}] (1.78,2.72) -- (1.78,0);
\node[below] at (0.44,0) {$1$};
\node[below] at (0.88,0) {$2$};
\node[below] at (1.34,0) {$3$};
\node[below] at (1.78,0) {$4$};
\end{tikzpicture}
\end{center}
}{
\smallskip\noindent\textbf{Решение от GPT}\par
Переходы 1 и 2 направлены вверх и соответствуют поглощению, причём $\nu_1=(E_1-E_0)/h<\nu_2=(E_2-E_0)/h$, поэтому А — 1. Переходы 3 и 4 направлены вниз и соответствуют излучению, причём $\varepsilon_3=E_3-E_0<\varepsilon_4=E_4-E_0$, поэтому Б — 3. Ответ: А — 1, Б — 3.
}

% Source 9F8E6C: https://ege.fipi.ru/bank/questions.php?search=1&proj=BA1F39653304A5B041B656915DC36B38&qid=9F8E6C
\BankTask{9F8E6C}{
На рисунке изображена упрощённая диаграмма нижних энергетических уровней атома. Нумерованными стрелками отмечены некоторые возможные переходы атома между этими уровнями. Какие из этих переходов связаны с поглощением кванта света наибольшей длины волны и излучением кванта света с наименьшей энергией?

Установите соответствие между процессами поглощения и испускания света и стрелками, обозначающими энергетические переходы атома.

\begin{center}
\TaskMatchingColumns[0.64\linewidth]{Процессы}{
А) поглощение кванта света наибольшей длины волны\par
Б) излучение кванта света с наименьшей энергией
}{Энергетические переходы}{
1) 1\par
2) 2\par
3) 3\par
4) 4
}
\end{center}

\begin{center}
\begin{tabular}{|c|c|}
\hline
А & Б \\
\hline
\phantom{1} & \phantom{3} \\
\hline
\end{tabular}
\end{center}

\begin{center}
\begin{tikzpicture}[book diagram, x=1cm, y=1cm]
\draw (0,-0.34) -- (0,3.62);
\draw[dashed] (0,3.34) -- (2.2,3.34);
\draw (0,2.78) -- (2.2,2.78);
\draw (0,2.50) -- (2.2,2.50);
\draw (0,1.94) -- (2.2,1.94);
\draw (0,1.10) -- (2.2,1.10);
\draw (0,0) -- (2.2,0);
\node[left] at (0,3.34) {$0$};
\node[left] at (0,2.78) {$E_4$};
\node[left] at (0,2.50) {$E_3$};
\node[left] at (0,1.94) {$E_2$};
\node[left] at (0,1.10) {$E_1$};
\node[left] at (0,0) {$E_0$};
\draw[-{book tip}] (0.46,0) -- (0.46,1.10);
\draw[-{book tip}] (0.92,0) -- (0.92,1.94);
\draw[-{book tip}] (1.38,2.50) -- (1.38,0);
\draw[-{book tip}] (1.84,2.78) -- (1.84,0);
\node[below] at (0.46,0) {$1$};
\node[below] at (0.92,0) {$2$};
\node[below] at (1.38,0) {$3$};
\node[below] at (1.84,0) {$4$};
\end{tikzpicture}
\end{center}
}{
\smallskip\noindent\textbf{Решение от GPT}\par
Энергия кванта равна разности энергий уровней: $\varepsilon=|E_f-E_i|=hc/\lambda$. Поглощению соответствуют стрелки 1 и 2; поскольку $E_1-E_0<E_2-E_0$, наибольшая длина волны соответствует стрелке 1. Излучению соответствуют стрелки 3 и 4; поскольку $E_3-E_0<E_4-E_0$, наименьшую энергию имеет квант при переходе 3. Ответ: А — 1, Б — 3.
}

% Source 667358: https://ege.fipi.ru/bank/questions.php?search=1&proj=BA1F39653304A5B041B656915DC36B38&qid=667358
\BankTask{667358}{
На рисунке изображена упрощённая диаграмма энергетических уровней атома. Нумерованными стрелками отмечены некоторые возможные переходы атома между этими уровнями. Какие из этих четырёх переходов связаны с поглощением света наибольшей длины волны и излучением света наибольшей длины волны?

Установите соответствие между процессами поглощения и испускания света и соответствующими переходами между уровнями энергии.

\begin{center}
\TaskMatchingColumns[0.64\linewidth]{Процесс}{
А) поглощение света наибольшей длины волны\par
Б) излучение света наибольшей длины волны
}{Переход}{
1) $E_0\to E_1$\par
2) $E_0\to E_2$\par
3) $E_3\to E_0$\par
4) $E_4\to E_0$
}
\end{center}

\begin{center}
\begin{tabular}{|c|c|}
\hline
А & Б \\
\hline
\phantom{1} & \phantom{3} \\
\hline
\end{tabular}
\end{center}

\begin{center}
\begin{tikzpicture}[book diagram, x=1cm, y=1cm]
\draw (0,-0.34) -- (0,3.68);
\draw[dashed] (0,3.38) -- (2.2,3.38);
\draw (0,2.82) -- (2.2,2.82);
\draw (0,2.54) -- (2.2,2.54);
\draw (0,1.98) -- (2.2,1.98);
\draw (0,1.14) -- (2.2,1.14);
\draw (0,0) -- (2.2,0);
\node[left] at (0,3.38) {$0$};
\node[left] at (0,2.82) {$E_4$};
\node[left] at (0,2.54) {$E_3$};
\node[left] at (0,1.98) {$E_2$};
\node[left] at (0,1.14) {$E_1$};
\node[left] at (0,0) {$E_0$};
\draw[-{book tip}] (0.46,0) -- (0.46,1.14);
\draw[-{book tip}] (0.94,0) -- (0.94,1.98);
\draw[-{book tip}] (1.40,2.54) -- (1.40,0);
\draw[-{book tip}] (1.88,2.82) -- (1.88,0);
\node[below] at (0.46,0) {$1$};
\node[below] at (0.94,0) {$2$};
\node[below] at (1.40,0) {$3$};
\node[below] at (1.88,0) {$4$};
\end{tikzpicture}
\end{center}
}{
\smallskip\noindent\textbf{Решение от GPT}\par
Длина волны $\lambda=hc/|E_f-E_i|$ максимальна при минимальной разности энергий уровней. Переходы 1 и 2 соответствуют поглощению, причём $E_1-E_0<E_2-E_0$, поэтому для А подходит переход 1. Переходы 3 и 4 соответствуют излучению, причём $E_3-E_0<E_4-E_0$, поэтому для Б подходит переход 3. Ответ: А — 1, Б — 3.
}

% Source 9AC530: https://ege.fipi.ru/bank/questions.php?search=1&proj=BA1F39653304A5B041B656915DC36B38&qid=9AC530
\BankTask{9AC530}{
На рисунке изображена упрощённая диаграмма энергетических уровней атома. Нумерованными стрелками отмечены некоторые возможные переходы атома между этими уровнями. Какие из этих четырёх переходов связаны с поглощением света наибольшей длины волны и излучением света с наименьшей длиной волны?

Установите соответствие между процессами поглощения и испускания света и соответствующими переходами между уровнями энергии.

\begin{center}
\TaskMatchingColumns[0.64\linewidth]{Процесс}{
А) поглощение света наибольшей длины волны\par
Б) излучение света наименьшей длины волны
}{Переход}{
1) $E_0\to E_1$\par
2) $E_0\to E_2$\par
3) $E_3\to E_0$\par
4) $E_4\to E_0$
}
\end{center}

\begin{center}
\begin{tabular}{|c|c|}
\hline
А & Б \\
\hline
\phantom{1} & \phantom{4} \\
\hline
\end{tabular}
\end{center}

\begin{center}
\begin{tikzpicture}[book diagram, x=1cm, y=1cm]
\draw (0,-0.34) -- (0,3.68);
\draw[dashed] (0,3.38) -- (2.2,3.38);
\draw (0,2.82) -- (2.2,2.82);
\draw (0,2.54) -- (2.2,2.54);
\draw (0,1.98) -- (2.2,1.98);
\draw (0,1.14) -- (2.2,1.14);
\draw (0,0) -- (2.2,0);
\node[left] at (0,3.38) {$0$};
\node[left] at (0,2.82) {$E_4$};
\node[left] at (0,2.54) {$E_3$};
\node[left] at (0,1.98) {$E_2$};
\node[left] at (0,1.14) {$E_1$};
\node[left] at (0,0) {$E_0$};
\draw[-{book tip}] (0.46,0) -- (0.46,1.14);
\draw[-{book tip}] (0.94,0) -- (0.94,1.98);
\draw[-{book tip}] (1.40,2.54) -- (1.40,0);
\draw[-{book tip}] (1.88,2.82) -- (1.88,0);
\node[below] at (0.46,0) {$1$};
\node[below] at (0.94,0) {$2$};
\node[below] at (1.40,0) {$3$};
\node[below] at (1.88,0) {$4$};
\end{tikzpicture}
\end{center}
}{
\smallskip\noindent\textbf{Решение от GPT}\par
Для кванта света $\lambda=hc/|E_f-E_i|$. Поглощению соответствуют переходы 1 и 2; меньшая разность $E_1-E_0<E_2-E_0$ даёт наибольшую длину волны, поэтому А — 1. Излучению соответствуют переходы 3 и 4; большая разность $E_4-E_0>E_3-E_0$ даёт наименьшую длину волны, поэтому Б — 4. Ответ: А — 1, Б — 4.
}

% Source 6A3D3F: https://ege.fipi.ru/bank/questions.php?search=1&proj=BA1F39653304A5B041B656915DC36B38&qid=6A3D3F
\BankTask{6A3D3F}{
На рисунке изображена упрощённая диаграмма нижних энергетических уровней атома. Нумерованными стрелками отмечены некоторые возможные переходы атома между этими уровнями. Какие из этих переходов связаны с поглощением кванта света наибольшей частоты и излучением кванта света наименьшей частоты?

Установите соответствие между процессами поглощения и излучения света и энергетическими переходами атома, указанными стрелками.

\begin{center}
\TaskMatchingColumns[0.58\linewidth]{Процессы}{
А) поглощение кванта света наибольшей частоты\par
Б) излучение кванта света наименьшей частоты
}{Энергетические переходы}{
1) 1\par
2) 2\par
3) 3\par
4) 4
}

\medskip
\begin{tabular}{|c|c|}
\hline
А & Б \\
\hline
\phantom{2} & \phantom{3} \\
\hline
\end{tabular}
\end{center}

\begin{center}
\begin{tikzpicture}[book diagram,x=0.035cm,y=0.035cm]
\draw (0,-5)--(0,181);
\draw (0,0)--(110,0);
\draw (0,55)--(110,55);
\draw (0,97)--(110,97);
\draw (0,125)--(110,125);
\draw (0,139)--(110,139);
\draw[dash pattern=on 4pt off 4pt] (7,167)--(110,167);
\node[left] at (0,0) {$E_0$};
\node[left] at (0,55) {$E_1$};
\node[left] at (0,97) {$E_2$};
\node[left] at (0,125) {$E_3$};
\node[left] at (0,145) {$E_4$};
\node[left] at (0,167) {$0$};
\draw[-{book tip}] (23,0)--(23,55);
\draw[-{book tip}] (46,0)--(46,97);
\draw[-{book tip}] (69,125)--(69,0);
\draw[-{book tip}] (92,139)--(92,0);
\node[below] at (23,0) {$1$};
\node[below] at (46,0) {$2$};
\node[below] at (69,0) {$3$};
\node[below] at (92,0) {$4$};
\end{tikzpicture}
\end{center}
}{
\smallskip\noindent\textbf{Решение от GPT}\par
Частота кванта определяется разностью энергий уровней: $\nu=|E_f-E_i|/h$. Переходы 1 и 2 соответствуют поглощению; поскольку $E_2-E_0>E_1-E_0$, наибольшая частота соответствует переходу 2. Переходы 3 и 4 соответствуют излучению; поскольку $E_3-E_0<E_4-E_0$, наименьшая частота соответствует переходу 3. Ответ: А — 2, Б — 3.
}

% Source 8DC935: https://ege.fipi.ru/bank/questions.php?search=1&proj=BA1F39653304A5B041B656915DC36B38&qid=8DC935
\BankTask{8DC935}{
На рисунке изображена упрощённая диаграмма нижних энергетических уровней атома. Нумерованными стрелками отмечены некоторые возможные переходы атома между этими уровнями. Какой из этих четырёх переходов связан с поглощением кванта света наименьшей частоты, а какой – с излучением кванта света наибольшей частоты?

Установите соответствие между процессами поглощения и излучения света и энергетическими переходами атома, указанными стрелками.

\begin{center}
\TaskMatchingColumns[0.58\linewidth]{Процессы}{
А) поглощение кванта света наименьшей частоты\par
Б) излучение кванта света наибольшей частоты
}{Энергетические переходы}{
1) 1\par
2) 2\par
3) 3\par
4) 4
}

\medskip
\begin{tabular}{|c|c|}
\hline
А & Б \\
\hline
\phantom{3} & \phantom{2} \\
\hline
\end{tabular}
\end{center}

\begin{center}
\begin{tikzpicture}[book diagram,x=0.035cm,y=0.035cm]
\draw (0,-5)--(0,139);
\draw (0,0)--(125,0);
\draw (0,45)--(125,45);
\draw (0,74)--(125,74);
\draw (0,90)--(125,90);
\draw (0,106)--(125,106);
\draw[dash pattern=on 24pt off 24pt] (0,125)--(125,125);
\node[left] at (0,0) {$E_0$};
\node[left] at (0,45) {$E_1$};
\node[left] at (0,74) {$E_2$};
\node[left] at (0,90) {$E_3$};
\node[left] at (0,106) {$E_4$};
\node[left] at (0,125) {$0$};
\draw[-{book tip}] (19,45)--(19,0);
\draw[-{book tip}] (44,74)--(44,0);
\draw[-{book tip}] (70,0)--(70,90);
\draw[-{book tip}] (94,0)--(94,106);
\node[below] at (19,0) {$1$};
\node[below] at (44,0) {$2$};
\node[below] at (70,0) {$3$};
\node[below] at (94,0) {$4$};
\end{tikzpicture}
\end{center}
}{
\smallskip\noindent\textbf{Решение от GPT}\par
Частота кванта определяется разностью энергий уровней: $\nu=|E_f-E_i|/h$. Переходы 3 и 4 соответствуют поглощению; поскольку $E_3-E_0<E_4-E_0$, наименьшая частота соответствует переходу 3. Переходы 1 и 2 соответствуют излучению; поскольку $E_2-E_0>E_1-E_0$, наибольшая частота соответствует переходу 2. Ответ: А — 3, Б — 2.
}

% Source F37670: https://ege.fipi.ru/bank/questions.php?search=1&proj=BA1F39653304A5B041B656915DC36B38&qid=F37670
\BankTask{F37670}{
На рисунке изображена упрощённая диаграмма нижних энергетических уровней атома. Нумерованными стрелками отмечены некоторые возможные переходы атома между этими уровнями. Какой из этих четырёх переходов связан с поглощением света наименьшей частоты, а какой — с излучением света наибольшей частоты?

Установите соответствие между процессами поглощения и испускания света и стрелками, указывающими энергетические переходы атома.

\begin{center}
\TaskMatchingColumns[0.60\linewidth]{Процессы}{
А) поглощение света наименьшей частоты\par
Б) излучение света наибольшей частоты
}{Энергетические переходы}{
1) 1\par
2) 2\par
3) 3\par
4) 4
}
\end{center}

\begin{center}
\begin{tabular}{|c|c|}
\hline
А & Б \\
\hline
\phantom{3} & \phantom{2} \\
\hline
\end{tabular}
\end{center}

\begin{center}
\begin{tikzpicture}[book diagram,x=0.035cm,y=0.035cm]
\draw (0,-10) -- (0,183);
\draw[dash pattern=on 4pt off 4pt] (7,169) -- (112,169);
\node[left] at (0,169) {$0$};
\draw (0,141) -- (112,141);
\node[left] at (0,141) {$E_4$};
\draw (0,127) -- (112,127);
\node[left] at (0,127) {$E_3$};
\draw (0,98) -- (112,98);
\node[left] at (0,98) {$E_2$};
\draw (0,56) -- (112,56);
\node[left] at (0,56) {$E_1$};
\draw (0,0) -- (112,0);
\node[left] at (0,0) {$E_0$};
\draw[-{book tip}] (23,56) -- (23,0);
\draw[-{book tip}] (46,98) -- (46,0);
\draw[-{book tip}] (70,0) -- (70,127);
\draw[-{book tip}] (93,0) -- (93,141);
\node[below] at (23,0) {$1$};
\node[below] at (46,0) {$2$};
\node[below] at (70,0) {$3$};
\node[below] at (93,0) {$4$};
\end{tikzpicture}
\end{center}
}{
\smallskip\noindent\textbf{Решение от GPT}\par
Частота фотона определяется соотношением $h\nu=|E_f-E_i|$. Поглощению соответствуют стрелки 3 и 4, причём $E_3-E_0<E_4-E_0$, поэтому наименьшая частота соответствует переходу 3. Излучению соответствуют стрелки 1 и 2; поскольку $E_2-E_0>E_1-E_0$, наибольшая частота соответствует переходу 2. Ответ: А — 3, Б — 2.
}

% Source 90AF7F: https://ege.fipi.ru/bank/questions.php?search=1&proj=BA1F39653304A5B041B656915DC36B38&qid=90AF7F
\BankTask{90AF7F}{
На рисунке изображена упрощённая диаграмма нижних энергетических уровней атома. Стрелками отмечены некоторые возможные переходы атома между этими уровнями.

Установите соответствие между процессами поглощения фотона наименьшей длины волны и излучения фотона наименьшей частоты и энергией соответствующего фотона.

\begin{center}
\TaskMatchingColumns[0.60\linewidth]{Процесс}{
А) поглощение фотона наименьшей длины волны\par
Б) излучение фотона наименьшей частоты
}{Энергия фотона}{
1) $E_1-E_0$\par
2) $E_2-E_0$\par
3) $E_3-E_0$\par
4) $E_4-E_0$
}
\end{center}

\begin{center}
\begin{tabular}{|c|c|}
\hline
А & Б \\
\hline
\phantom{3} & \phantom{1} \\
\hline
\end{tabular}
\end{center}

\begin{center}
\begin{tikzpicture}[book diagram,x=0.035cm,y=0.035cm]
\draw[gray] (0,-28) -- (0,140);
\draw[densely dashed] (4,135) -- (95,135);
\node[left] at (0,135) {$0$};
\draw (2,111) -- (83,111);
\node[left] at (0,111) {$E_4$};
\draw (2,95) -- (83,95);
\node[left] at (0,95) {$E_3$};
\draw (2,69) -- (83,69);
\node[left] at (0,69) {$E_2$};
\draw (2,39) -- (83,39);
\node[left] at (0,39) {$E_1$};
\draw[gray] (-20,0) -- (99,0);
\node[below left] at (0,0) {$E_0$};
\draw[-{book tip}] (23,39) -- (23,0);
\draw[-{book tip}] (41,111) -- (41,0);
\draw[-{book tip}] (55,0) -- (55,69);
\draw[-{book tip}] (73,0) -- (73,95);
\end{tikzpicture}
\end{center}
}{
\smallskip\noindent\textbf{Решение от GPT}\par
Энергия фотона равна $\varepsilon=h\nu=hc/\lambda=|E_f-E_i|$. Из показанных переходов поглощения $E_0\to E_2$ и $E_0\to E_3$ наименьшей длине волны соответствует большая энергия $E_3-E_0$ — вариант 3. Из переходов излучения $E_1\to E_0$ и $E_4\to E_0$ наименьшей частоте соответствует меньшая энергия $E_1-E_0$ — вариант 1. Ответ: А — 3, Б — 1.
}

\end{document}
